% !TeX root = ./main.tex

\sysusetup{
  title       = {课程论文标题},
  course-paper = true,
  author      = {作者姓名},
  department  = {院系名称},
  student-id  = {学号},
  cover-heading = {课程论文},
  header-title  = {课程论文},
  cover-items = {
    \sysucoverrow{课程名称}{课程名称}
    \sysucoverrow{姓名}{作者姓名}
    \sysucoverrow{学号}{学号}
    \sysucoverrow{院系}{院系名称}
    \sysucoverrow{任课教师}{教师姓名}
  },
}

% 自动检测我们定制的 Dev Container 环境并替换 CJK 字体，保持跨平台兼容性
\IfFileExists{/.sysu-thesis.devcontainer}{
  \sysusetup{cjk-font = noto}
}{}

%\renewcommand{\thefigure}{\thechapter-\arabic{figure}}% 将正文和目录的图片标题由x.x改为x-x，慎用
%\renewcommand{\thetable}{\arabic{chapter}-\arabic{table}}% 将正文和目录的表格标题由x.x改为x-x，慎用
%\renewcommand{\theequation}{\arabic{chapter}-\arabic{equation}}% 将正文的公式序号由(x.x)改为(x-x)，慎用

% 加载宏包

% 定理类环境宏包
\usepackage{amsthm}

% 中英文双标题
\usepackage{bicaption}

% 表注
\usepackage{threeparttable}

% 跨页表格
\usepackage{longtable}

% 算法
\usepackage[ruled,linesnumbered]{algorithm2e}
%\renewcommand{\thealgocf}{\arabic{chapter}-\arabic{algocf}}% 将算法标题由x.x改为x-x，慎用

% 代码块
\usepackage{listings}

% SI 量和单位
\usepackage{siunitx}

% 参考文献使用 BibTeX + natbib 宏包
% 顺序编码制
\usepackage[sort]{natbib}
\bibliographystyle{sysuthesis-numerical}
% 著者-出版年制
% \usepackage{natbib}
% \bibliographystyle{sysuthesis-authoryear}

% 参考文献使用 BibLaTeX 宏包，下面三行代码三选一。
% \usepackage[style=sysuthesis-numeric]{biblatex}
% \usepackage[bibstyle=sysuthesis-numeric,citestyle=sysuthesis-inline]{biblatex}
% \usepackage[style=sysuthesis-authoryear]{biblatex}
% 声明 BibLaTeX 的数据库
% \addbibresource{reference.bib}

% 配置图片的默认目录
% \graphicspath{{image/}}

% 数学命令
\makeatletter
\newcommand\dif{%  % 微分符号
  \mathop{}\!%
  \ifsysu@math@style@TeX
    d%
  \else
    \mathrm{d}%
  \fi
}
\makeatother
\newcommand\eu{{\symup{e}}}
\newcommand\iu{{\symup{i}}}

% 用于写文档的命令
\DeclareRobustCommand\cs[1]{\texttt{\char`\\#1}}
\DeclareRobustCommand\env[1]{\texttt{#1}}
\DeclareRobustCommand\pkg[1]{\textsf{#1}}
\DeclareRobustCommand\file[1]{\nolinkurl{#1}}

% hyperref 宏包在最后调用
\usepackage{hyperref}
